% titel: Blatt 0, A2: Binomische Formeln

% Teil I: Werkzeuge, kein Lösungsweg.

Ausmultiplizieren von Termen, wobei $x$, $y$ Variablen sind: die drei
binomischen Formeln.

\begin{werkzeugkasten}[Binomische Formeln]
\begin{align}
  (x+y)^2 &= x^2+2xy+y^2, \\
  (x-y)^2 &= x^2-2xy+y^2, \\
  (x+y)(x-y) &= x^2-y^2 .
\end{align}
Rechts steht bei der ersten Formel das \emph{doppelte} Produkt $2xy$; der
Mittelterm geht beim Quadrieren einer Differenz mit Minuszeichen ein.
\end{werkzeugkasten}

\begin{werkzeugkasten}[Ausmultiplizieren allgemein]
Ohne binomische Formel wendet man das Distributivgesetz an: jeden Summanden
der ersten Klammer mit jedem der zweiten multiplizieren und die Teilstücke
addieren,
\begin{equation}
  (a+b)(c+d)=ac+ad+bc+bd .
\end{equation}
Gleichnamige Terme zusammenfassen.
\end{werkzeugkasten}

\begin{achtung}[Darauf achten]
\begin{itemize}
  \item $(x+y)^2\neq x^2+y^2$ -- der Term $2xy$ darf nicht fehlen.
  \item Bei $(x-y)^2$ ist der Mittelterm $-2xy$, nicht $+2xy$.
  \item $(x+y)(x-y)=x^2-y^2$: das mittlere $(+xy-xy)$ kürzt sich weg.
  \item Faktoren wie $2$ oder Vorzeichen aus der Klammer korrekt mittels
        Distributivgesetz verteilen.
\end{itemize}
\end{achtung}
